\documentclass[11pt,a4paper]{article}
\usepackage[T1]{fontenc}
\usepackage[utf8]{inputenc}
\usepackage{lmodern,amsmath,amssymb,amsthm,mathtools}
\usepackage[margin=24mm]{geometry}
\usepackage{microtype,booktabs,enumitem}
\usepackage[colorlinks=true,linkcolor=blue,urlcolor=blue]{hyperref}
\newtheorem{lemma}{Lemma}
\newcommand{\Z}{\mathbb Z}
\newcommand{\e}{\mathbf e}
\newcommand{\pos}[1]{[#1]_+}
\newcommand{\strip}{\mathrm{strip}}
\setlength{\parindent}{0pt}
\setlength{\parskip}{5pt}
\setlist{nosep,leftmargin=*}
\title{Deleting the middle edge of an FG strip\\\large Proof blueprint for a theorem of \emph{Ice and mirrors}}
\author{Cluster HMS}
\date{14 November 2026}
\begin{document}
\maketitle
\vspace{-1.5em}
This is the supplied mathematical proof to formalize, from \texttt{main.tex}.
The source commit is
\begin{center}\texttt{7e310e01df64e38be583500fc81a8ef1a001a196}.\end{center}
The companion \emph{Target theorem} specifies the universal Lean statement.
The identities below are proof obligations, not reports of completed Lean proofs.
The choice of formal definitions and proof decomposition is left to the formalization.

\section{The matrix calculation and its ambient scope}
For every $n\ge2$ and $N\ge\max\{\lfloor n/2\rfloor,2\}$, use the strip
\[
 \tau_{2i+1}=(T_{i+1},T_i,B_i),\quad
 \tau_{2i+2}=(B_i,B_{i+1},T_{i+1})\quad(0\le i<N),
\]
with $e_{2i}=T_iB_i$, $e_{2i+1}=T_{i+1}B_i$ and distinguished edge $e_N$.
The lattice points in each triangle have $a+b+c=n$ and exclude the corners.
The signed matrix is the sum of the cycles
\begin{equation}\label{eq:cycles}
 (a+1,b+1,c)\longrightarrow(a,b+1,c+1)
 \longrightarrow(a+1,b,c+1)\longrightarrow(a+1,b+1,c),
 \quad a+b+c=n-2.
\end{equation}
Glue equal edge weights and omit the actual top and bottom boundary vertices.
Positive $B_{uv}$ denotes $u\to v$. Put $\rho_Q(v)=B_{v,*}$.
The vector $\e_w$ is the coordinate vector on the current live set;
it is zero when the label $w$ is absent. Repeated labels contribute with
their summed coefficients.

Write $I=V_{\strip}\setminus(L(e_0)\cup L(e_{2N}))$.
The ambient matrix on a finite set $V$ is
\[
 B=A+E,\qquad A=B_{\strip,\mathrm{ext}},\qquad
 E^\mathsf{T}=-E,\qquad E_{\iota(x),v}=0\ (x\in I,\ v\in V).
\]
If $E_{k,*}=0$, skew-symmetry gives $E_{*,k}=0$. Substitution in matrix mutation gives
\begin{equation}\label{eq:ambient-mutation}
 \mu_k(A+E)=\mu_k(A)+E.
\end{equation}
Thus a word supported in $\iota(I)$ has the same complete pivot and deletion
vertex expressions as the strip computation, with zero exterior entries.
The summand $E$ remains unchanged throughout.
This also permits arbitrary attachments along the unmutated sides of a
local triangle or two-triangle configuration.

The proof removes triangle interiors and then successive edge layers.
The local words avoid the interface edges of their configurations.
Mutation commutes with restriction to any live subset containing the pivot.
Together with $\mu_k^2=\mathrm{id}$ this gives, for $s\notin W$,
\begin{equation}\label{eq:undo}
 \mu_{W^{\mathrm{rev}}}\bigl((\mu_W B)|_{U\setminus\{s\}}\bigr)
 =B|_{U\setminus\{s\}}.
\end{equation}
Consequently each completed deletion move restores every retained entry.
Additional deletions away from its word and its deletion vertex preserve
both its calculation and its source or sink condition.

\newpage
\section{Deleting triangle interiors}
Number an edge by $1\le k\le n-1$. Deletion to depth $p$ removes
$k\le p$ or $k\ge n-p$. In a triangle with boundary side $c=0$, put
\begin{equation}\label{eq:Cp}
 C_p=\{(a,b,c)\in\Z_{>0}^3:a+b+c=n,\quad
                   a+c\le p\ \text{or}\ b+c\le p\}.
\end{equation}
Thus $C_0=C_1=\varnothing$. The retained sides before deletions are
$D_k=(n-k,0,k)$ and $R_k=(0,n-k,k)$.

\begin{lemma}\label{lem:triangle}
Suppose $1\le p\le(n-1)/2$. Delete both nonboundary edges to depth $p$
and delete $C_p$. There are deletion moves removing $C_{p+1}\setminus C_p$
whose words and deleted vertices are all interior to the triangle.
\end{lemma}

The source proof deletes $a+c=p+1$ in decreasing order of $a$, then exchanges
$a$ and $b$. If $n=3$, $p=1$, the only interior vertex is isolated.
For $n\ge4$, put $b_0=n-p-1$. At the deletion of
\[
 s_a=(a,b_0,p+1-a)\qquad(a=p,p-1,\ldots,1),
\]
use the word
\begin{equation}\label{eq:triangle-word}
 v_1,\ldots,v_{b_0-1},\qquad v_b=(a,b,n-a-b),\qquad v_{b_0}=s_a.
\end{equation}
Every pivot is interior and remains present, since
\[
 a+(n-a-b)=n-b\ge p+2,\qquad
 b+(n-a-b)=n-a\ge n-p\ge p+1.
\]
The word also avoids the already deleted $s_{a'}$ with $a'>a$.

Put $A_b=(a-1,b+1,n-a-b)$ and $B_b=(a+1,b,n-a-b-1)$.
The incoming neighbours of $v_1$ are the present vertices among $A_1,B_1$.
The neighbour $(a,0,n-a)$ is deleted from its edge. The exceptions are
\[
 A_1=(0,2,n-2)\text{ absent if }a=1,\ p\ge2;
 \qquad B_1\in C_p\text{ if }n=2p+1,\ a=p.
\]
For $2\le b\le b_0-1$, the complete expression before mutation is
\begin{equation}\label{eq:triangle-pivot}
 \rho_{\mathrm{pre}}(v_b)=\e_{v_{b-1}}+\e_{v_{b+1}}-\e_{A_b}-\e_{B_b}.
\end{equation}
The cycles in \eqref{eq:cycles} give this successively. Earlier pivots
$v_1,\ldots,v_{b-2}$ do not change the arrows at $v_b$.
Paths through $v_{b-1}$ cancel the outgoing arrows off the line
$a=\mathrm{constant}$, and $v_{b-1}\to v_b$ reverses.

Before the last mutation, the deletion vertex satisfies
\[
 \rho_{\mathrm{pre}}(s_a)=\e_{A_{b_0-1}}+\e_{B_{b_0-1}}-\e_{v_{b_0-1}}.
\]
Its other possible neighbours lie in $C_p$, on a deleted edge portion,
on $c=0$, or at $s_{a+1}$. Mutation at $v_{b_0-1}$ therefore gives
\begin{equation}\label{eq:triangle-final}
 \rho_{\mathrm{final}}(s_a)=\e_{v_{b_0-1}}.
\end{equation}
Delete $s_a$ and undo the word using \eqref{eq:undo}. Exchanging $a,b$
reverses all arrows and gives the other family of deletion moves.
Its pivots have $b\le p$, whereas the first deleted family has $b=b_0\ge p$.
When $b_0=p$, equivalently $n=2p+1$, skip $(p,p,1)$, already deleted.
These are the endpoint and overlap cases of Lemma~\ref{lem:triangle}.

\newpage
\section{Two triangles and the common-edge word}
Place the two oriented triangles in $[0,n]^2$ by
\[
 x(a,b,c)=(b+c,c),\qquad y(a,b,c)=(b,b+c).
\]
Their omitted boundary sides are $c=0$ in $x$ and $a=0$ in $y$.
The labels are
\begin{align*}
 D_k&=x(n-k,0,k)=y(n-k,k,0)=(k,k),&&1\le k\le n-1,\\
 X_{i,j}&=x(n-i,i-j,j)=(i,j),&&1\le j<i\le n-1,\\
 Y_{i,j}&=y(n-j,i,j-i)=(i,j),&&1\le i<j\le n-1,\\
 L_j&=y(n-j,0,j)=(0,j),&&1\le j\le n-1,\\
 R_j&=x(0,n-j,j)=(n,j),&&1\le j\le n-1.
\end{align*}
Set $X_{n,j}=R_j$. Boundary or deleted labels contribute zero.
For the $y$ triangle, use $(b,c,a)$ in \eqref{eq:Cp}.
The map
\begin{equation}\label{eq:halfturn}
 H(i,j)=(n-i,n-j)
\end{equation}
preserves the arrows and interchanges $D_k,D_{n-k}$.
Every adjacent pair of strip triangles has this matrix or its negative.
Both identifications must follow from the glued cycles, including all
vertices on $L$ and $R$.

\begin{lemma}\label{lem:edge}
Suppose $0\le p<\lfloor n/2\rfloor$. Delete the three edges $L,D,R$ to
depth $p$ and delete $C_{p+1}$ in each triangle. Deletion moves remove
$D_{n-p-1}$ and $D_{p+1}$, counted only once if they coincide.
All words and deleted vertices lie in the two triangle interiors or on $D$.
\end{lemma}

When $n=2p+2$, the two endpoints coincide and the source gives
\[
 \rho_Q(D_{p+1})=-\e_{x(p,1,p+1)}-\e_{y(p+1,p,1)}.
\]
It is a sink, with the absent-label convention; no mutation is needed.
This includes the rank-two calculation at $p=0$.

For $n\ge2p+3$, put
\[
 r=p+2,\qquad t=n-p-1,\qquad s=D_t.
\]
The word that makes $s$ a sink is
\begin{equation}\label{eq:edge-word}
 D_{p+1},\ldots,D_{t-1};\quad
 X_{r,r-1},\ldots,X_{r,1};\quad\cdots;\quad
 X_{t,t-1},\ldots,X_{t,1}.
\end{equation}
The semicolons only separate successive portions of the word.
For an $X$ vertex, membership in $C_{p+1}$ means
$i\le p+1$ or $i-j\ge t$. Every displayed $X_{i,j}$ remains present because
\[
 r\le i\le t,\qquad 1\le j<i,\qquad i-j\le t-1.
\]
All displayed diagonal vertices remain present as well.
The word avoids $s$ and every $L_k,R_k$. After it, delete $s$ and undo
the mutations. The second endpoint is treated on the following pages.

\newpage
\section{Complete expressions during the common-edge word}
For $r\le i<t$, the expression before mutating the top vertex is
\begin{equation}\label{eq:top}
 \rho_{\mathrm{pre}}(X_{i,i-1})
 =-\e_{D_{i-1}}+\e_{D_i}-\e_{X_{i+1,i-1}}+\e_{X_{i,i-2}}.
\end{equation}
At $i=t$ it is
\begin{equation}\label{eq:last}
 \begin{split}
 \rho_{\mathrm{pre}}(X_{t,t-1})={}&-\e_{D_{t-1}}-\e_s-\e_{X_{t+1,t-1}}\\
 &+\e_{X_{t+1,t}}+\e_{Y_{t-1,t}}+\e_{X_{t,t-2}}.
 \end{split}
\end{equation}
For $r\le i\le t$ and $1\le j<i-1$, put
\[
 K_{ij}=\begin{cases}D_{p+1},&i=r,\\X_{i-1,j},&i>r.\end{cases}
\]
Then the lower vertices satisfy
\begin{equation}\label{eq:lower}
 \begin{split}
 \rho_{\mathrm{pre}}(X_{i,j})={}&-\e_{K_{ij}}-\e_{X_{i+1,j}}
       +\e_{X_{i,j+1}}+\e_{X_{i,j-1}}-\mathbf1_{i=t}\e_s.
 \end{split}
\end{equation}
These are identities on the entire live set. A proof must establish that
every unlisted coefficient is zero, including coefficients at interface vertices.

The diagonal part of \eqref{eq:edge-word} meets $X_{i,i-1}$ only through
$D_{i-1},D_i$. Both are mutated for $i<t$; $D_t=s$ is not mutated.
No preceding $X_{h,j}$ with $h<i$ changes this top expression.
This gives \eqref{eq:top} and \eqref{eq:last} from matrix mutation.

After the diagonal mutations, the lower expression is
\[
 \begin{split}
 \rho_Q(X_{i,j})={}&\e_{X_{i+1,j+1}}-\e_{X_{i+1,j}}-\e_{X_{i,j+1}}\\
 &+\e_{X_{i,j-1}}+\e_{X_{i-1,j}}-\e_{X_{i-1,j-1}}.
 \end{split}
\]
For $i=r$, mutation at $X_{r,j+1}$ cancels the arrow to $X_{r+1,j+1}$
and creates the arrow from $D_{p+1}$.
For $i>r$, only the preceding mutations at $X_{i-1,j}$ and $X_{i,j+1}$
affect $X_{i,j}$. The former cancels the arrow from $X_{i-1,j-1}$
and creates an arrow to $X_{i-1,j+1}$, interpreted as $D_{i-1}$ when
$j=i-2$. The latter cancels the arrows to $X_{i-1,j+1}$ and $X_{i+1,j+1}$.
Reversing the arrows incident to these two pivots gives \eqref{eq:lower}.

Only the mutations at $X_{t,j}$ affect $s$. Before the top one,
\[
 \rho_{\mathrm{pre}}(s)=\e_{X_{t,t-1}}-\e_{X_{t+1,t}}-\e_{Y_{t-1,t}}.
\]
The possible neighbour $Y_{t,t+1}$ is on the boundary for $p=0$ and lies
in $C_{p+1}$ for $p\ge1$. Mutation at $X_{t,t-1}$ cancels the two negative
terms. After mutating $X_{t,j}$, the expression becomes
\[
 \rho_Q(s)=-\e_{X_{t,j}}+\e_{X_{t,j-1}}
 \qquad(j=t-1,t-2,\ldots,1).
\]
Since $X_{t,0}$ is absent, the word ends with
\begin{equation}\label{eq:edge-final}
 \rho_{\mathrm{final}}(s)=-\e_{X_{t,1}}.
\end{equation}
This is the source proof's decisive sink calculation.

\newpage
\section{The second endpoint and additional deletions}
After deleting $D_t$ and undoing its word, apply the half-turn
\eqref{eq:halfturn}. In the rotated coordinates, $D_{p+1}$ is already
absent and the new deletion vertex is $D_t$. Hence the word in
\eqref{eq:edge-word} must omit its first mutation at $D_{p+1}$.
In the original coordinates the deleted endpoint was $D_t$, and the
new deletion vertex is $D_{p+1}$.

If $r<t$, replace $D_{p+1}$ by $Y_{p+1,p+2}$ in \eqref{eq:top}
and \eqref{eq:lower} at $i=r$. Explicitly, the modified top expression is
\[
 \rho_{\mathrm{pre}}(X_{r,r-1})
 =-\e_{Y_{p+1,p+2}}+\e_{D_r}
   -\e_{X_{r+1,r-1}}+\e_{X_{r,r-2}},
\]
and the lower expressions use $K_{rj}=Y_{p+1,p+2}$.
The mutation at $D_r$ leaves the arrow from $Y_{p+1,p+2}$ which the
omitted mutation would have cancelled. All other formulas, including
\eqref{eq:edge-final}, are unchanged.

If $r=t$, equivalently $n=2p+3$, the diagonal portion of the second word
is empty. Omit $-\e_{D_{t-1}}+\e_{Y_{t-1,t}}$ from \eqref{eq:last}, giving
\[
 \rho_{\mathrm{pre}}(X_{t,t-1})
 =-\e_s-\e_{X_{t+1,t-1}}+\e_{X_{t+1,t}}+\e_{X_{t,t-2}}.
\]
Omit $-\e_{K_{tj}}$ from \eqref{eq:lower} as well. In the rotated
coordinates the final expression is
\begin{equation}\label{eq:second-final}
 \rho_{\mathrm{final}}(D_t)=-\e_{X_{t,1}}-\e_{Y_{t-1,t}}.
\end{equation}
Both coefficients are nonpositive. Thus the second deletion is valid
also at this shortest nontrivial word. The coincident-endpoint case
$n=2p+2$ was handled before forming either word.

The modified word is not justified merely by restricting the first
calculation: the removed vertex was itself the first pivot. The modified
identities above must be proved separately from the mutation formula.
They are also required for $p=0$, including the rank-three case.

Other previously deleted interface vertices are different. No $L_j$ or
$R_j$ occurs as a pivot or deletion vertex in these words. For any subset
$F$ of those vertices, restriction commutes with each mutation, so
\[
 (\mu_W Q)|_{U\setminus F}=\mu_W(Q|_{U\setminus F}).
\]
Restricting a nonpositive or nonnegative expression preserves its sign.
This proves validity when neighbouring edge deletions have removed more
than the local depth hypothesis requires. Equation~\eqref{eq:ambient-mutation}
gives the same conclusion in an ambient quiver attached along $L,R$.

Arrow reversal is also harmless:
\[
 \mu_k(-B)=-\mu_k(B).
\]
It interchanges the source and sink conditions. These facts apply to the
alternating orientations of adjacent strip pairs. They do not permit
omitting a live interface coefficient when checking a complete expression.

\newpage
\section{The rounds for arbitrary strip length}
Put $q=\lfloor n/2\rfloor$. Apply Lemma~\ref{lem:edge} with $p=0$ to
$e_1,\ldots,e_{2N-1}$ in increasing order. Since $C_1=\varnothing$,
each common edge is deleted to depth $1$. Earlier moves only add
deletions on the next pair's $L$ or $R$ edge. The restriction argument
above proves that its words remain valid.

After round $p$, the source proof records
\begin{equation}\label{eq:round}
 \begin{aligned}
 &k\le p\text{ or }k\ge n-p\text{ deleted on }e_j,
       &&p\le j\le2N-p,\\
 &C_p\text{ deleted in }\tau_i,
       &&p\le i\le2N-p+1.
 \end{aligned}
\end{equation}
The completed matrix is the restriction of the original ambient matrix
to the remaining vertices. Every previous mutation and deletion belongs
to $\iota(I)$. Both assertions must be retained in the induction.

Suppose $1\le p<q$. For $p+1\le i\le2N-p$, both nonboundary edges of
$\tau_i$ have depth $p$. Apply Lemma~\ref{lem:triangle} to remove
$C_{p+1}\setminus C_p$ in each of these triangles. Then apply
Lemma~\ref{lem:edge} to
\[
 e_{p+1},e_{p+2},\ldots,e_{2N-p-1}
\]
in increasing order. Each adjacent pair has the required interior
deletions. Additional deletions on its exterior edges are allowed by
restriction. This proves \eqref{eq:round} for $p+1$.

More explicitly, the listed operations leave exactly depth
$\min\{p,j,2N-j\}$ on $e_j$, and delete exactly
$C_{\min\{p,i,2N-i+1\}}$ in $\tau_i$ after round $p\ge1$.
These equalities record the deletions outside the central range in
\eqref{eq:round}; no larger rank-dependent strip-length assumption is used.

Since $N\ge q$, the edge $e_N$ is in the range $q\le j\le2N-q$.
Its remaining vertices would satisfy $q<k<n-q$. If $n=2q$ this interval
is empty; if $n=2q+1$ it contains no integer. Thus every lattice vertex
on $e_N$ has been deleted. This argument uses the edge index $N$ and
does not assume $N$ is even.

No common edge in these rounds is $e_0$ or $e_{2N}$.
The triangle words are interior; the common-edge words avoid $L$ and $R$.
Hence every mutation and deletion belongs to $\iota(I)$.
Composing \eqref{eq:undo} gives
\[
 \iota(L(e_N))\subseteq Z\subseteq\iota(I),\qquad
 B_{\mathrm{final}}=B|_{V\setminus Z}.
\]

\section{Formal correspondence and verification}
Prove the equivalence between the chosen finite representation and the
glued triangle lattice before applying the local calculations to the
strip. For an embedded strip in a surface, split the triangle sum into
strip and exterior triangles. Every exterior summand has zero entries
at vertices in $\iota(I)$, which verifies the ambient hypothesis.

The source locators are \texttt{lem:triangle-shell} (lines 3575--3654),
\texttt{lem:edge-layer} (3860--3980), and the main theorem
(statement 3402--3409, proof 4119--4161). The FG construction is at
3134--3151; the two-triangle coordinates are at 3767--3792.
The delivered source-to-Lean map must name the final theorem, the local
results actually used, and all hypotheses of the ambient application.

Build in a pinned Lean/mathlib environment and check a fresh importing
consumer. The transitive \texttt{\#print axioms} audit permits only
\texttt{propext}, \texttt{Classical.choice}, and \texttt{Quot.sound}.
Finite exact computations are supplementary. The supplied checker passed
for $2\le n\le8$ during preparation. It uses
$N=n\mathbin{//}2$, has no arbitrary exterior attachment, and gives
$N=1$ for $n=2,3$, outside the paper's $N\ge2$ hypothesis.
It does not prove any universally quantified local identity or the
arbitrary-$N$ induction. Report any failed identity with its exact
source location, and distinguish unfinished obligations from a completed
kernel-checked theorem.
\end{document}
